\section{Sources and reproducibility}
\label{app:artifacts}

Relative paths below resolve from the repository root immediately above
\path{revision-2/}. Byte counts and SHA-256 digests for all 118 pinned
inputs appear in \path{revision-2/data/source_manifest.json}, paired with a
flat lookup table in \path{revision-2/data/source-index.csv}. Each entry
records the exact file consumed by our secondary-analysis script while
preserving the original campaign manifests as separate provenance records.

\begin{description}
\item[\hypertarget{artifact-E1}{Decoy-gradient experiment.}]
Primary text and tables reside in
\path{evidence/decoy-gradient/paper1-main.tex}, which matches public
preprint \cite{paper1} on the 4,096-frame support leak, the $+0.65$ to
$+1.50\pp$ token-probe gain, and the $0.01$-nat clipping-and-noise utility
penalty. Frame tallies remain source-reported here without re-parsing raw
training logs.

\item[\hypertarget{artifact-B1}{Boundary-replacement audit.}]
Archived findings live in
\path{evidence/boundary-replacement/independent_evidence_review_2026-09-09/VERDICTS.md}
and its companion recomputed JSON file. Those records trace the shifted
sequence alignment $(x_1,\ldots,x_{T-1})$ versus $(x_2,\ldots,x_T)$, the
six zero-return baseline comparisons, and the $16.64\pp$ answer-token drop
to their original evaluation outputs, which our manifest hashes directly.

\item[\hypertarget{artifact-R1}{Privacy, utility and gain exploration.}]
Raw exploratory outputs come from
\path{reproduction/trilemma-campaign/evidence/r1-exploration-20260917/analysis/exploratory_analysis.json},
\path{reproduction/trilemma-campaign/evidence/r1-qualification-20260916/attack-spec-frozen.json},
\path{evidence/trilemma-campaign/r1_cell_results.csv}, and
\path{evidence/trilemma-campaign/r1_seed_results.csv}. Our script rebuilds
all 76 configuration-by-seed contrasts, Student-$t$ endpoints, probe
qualifications, and gate outcomes, while reading the unexecuted planning
bound from
\path{reproduction/trilemma-campaign/evidence/r1-exploration-20260917/confirmation_power.json}.

\item[\hypertarget{artifact-P4}{Defence-redesign pilot.}]
Pilot records comprise
\path{reproduction/defense-redesign-pilot/outputs/runs/pilot-20260922-a/verification.json},
adjacent run-metrics files, and
\path{evidence/defense-redesign-pilot/PILOT_DECISION.md}. From those saved
candidate and null arrays, our script recomputes three-seed paired utility
bounds, atomic structural advantages, and all-bindings composite verdicts.

\item[\hypertarget{artifact-E8}{Orthogonal-wrapper experiment.}]
Wrapper measurements come from
\path{reproduction/structural-e8/paper-data/collected/h100-2/issues/issue_e8_27b_20260807/e8_robustness_27b_bf16_densek_20260807T230529Z.json}
and \path{evidence/structural-e8/EXPERIMENT_STORY.md}. Extracted fields
cover undefended baselines, depth-specific $K_{50}$ interpolation brackets,
pooled $3.71\%$ recovery, and individual session scores up to $5.73\%$.

\item[\hypertarget{artifact-ER}{Rotation deployment study.}]
Deployment records draw from
\path{reproduction/structural-e8/paper-data/collected/toula/er_20260811/COLLECTION_MANIFEST.json},
accompanying session-identity JSONs, and
\path{evidence/structural-e8/EVIDENCE_EXECUTION_LEDGER.md}. Across eight runs and 24 sessions, session~1 at \texttt{N=16} and zero RTT
logs $2.8$ versus $7.6$ tokens/s and establishes the $63.16\%$ throughput
drop at recorded precision.

\item[\hypertarget{artifact-MI}{MI-estimator analysis.}]
Variational critic records reside in
\path{evidence/privacy-contract/privacy-contract-draft.tex} and
\path{evidence/privacy-contract/KEY_NUMBERS.md}, paired with finite
architecture sweeps in
\path{evidence/structural-e8/LATENT_CAMPAIGN_REPORT_V7_V9.md}. Near-zero
critic scores are cited as source-reported alongside the analytic
counterexample in Eq.~\eqref{eq:dv}.
\end{description}

\section{Checking the derived records}
\label{app:checks}

Running the standalone Python script inside \path{revision-2/} regenerates
every derived table and consistency check.
\begin{quote}\small\ttfamily
python3 analysis/recompute.py
\end{quote}
Execution relies exclusively on the Python standard library and writes
outputs under \path{revision-2/data/} without touching model checkpoints,
tensor captures, or cryptographic keys. Saved logs record 1,522 checks.
Across the 118 hashed inputs, those assertions compare recomputed seed
contrasts, confidence limits, and gate flags against the archived campaign
tables.

\begin{table}[h]
\centering\small
\begin{tabular}{@{}lr@{}}
\toprule
Rejection reason combination & Configurations \\
\midrule
Qualified privacy breach only & 10 \\
Gain failure and qualified privacy breach & 3 \\
Utility and gain failure & 3 \\
Utility failure only & 2 \\
Gain failure only & 1 \\
\midrule
Total & 19 \\
\bottomrule
\end{tabular}
\caption{Disjoint rejection-reason intersections across all 19 exploratory
configurations (\artifact{R1}). Partitioning the marginal bars in
Figure~\ref{fig:campaigns} confirms that all six configurations without a
qualified privacy breach fail utility or learning-gain gates independently,
as stored in \texttt{data/audit\_results.json}.}
\label{tab:intersections}
\end{table}

Derived JSON outputs preserve per-campaign confidence conventions, scoring
denominators, and operational caveats. Vector scripts verify inputs first.
Before rendering PDFs, the plotting code cross-checks its CSV tables against
both the exploratory seed logs and the defence-pilot verification report,
while \path{build.py} enforces byte-level preservation of the historical
\path{paper/} directory.
